\section{Why Most Levers Failed}
\label{sec:discussion}

Most of what we tried did not survive held-out tasks (\cref{fig:transfer}). Every structure built
for hard tasks lost its lead on held-out tasks; the acceptance checker alone stayed ahead, and it
shipped only as a gate for specification-style requests. The failures share five causes. We
report them because each cost weeks and would cost others the same.

\begin{figure}[t]
\centering
\begin{tikzpicture}
\begin{axis}[
  width=0.50\textwidth, height=5.4cm,
  xmin=-35, xmax=55, ymin=0.4, ymax=5.6, y dir=reverse,
  xtick={-20,0,20,40}, minor xtick={-30,-10,10,30,50}, xminorgrids,
  minor grid style={ruleGray!25},
  ytick={1,...,5},
  yticklabels={
    {Reading rules on every node (V1)},
    {Three builders and a judge (V1)},
    {Three readers and a reconciler (V1)},
    {Acceptance checker (V1)},
    {Check at first delivery, O2 (V8)}},
  yticklabel style={font=\footnotesize},
  tick label style={font=\footnotesize},
  xlabel={Arm minus same-batch base (percentage points)},
  xlabel style={font=\footnotesize},
  xmajorgrids, grid style={ruleGray!45},
  axis x line*=bottom, axis y line*=left,
  y axis line style={draw=none}, ytick style={draw=none},
  legend style={font=\footnotesize, draw=none, fill=none, at={(0.5,1.03)}, anchor=south,
    legend columns=2, /tikz/every even column/.append style={column sep=8pt}},
  clip=false,
]
  \draw[ink!55, dashed] (axis cs:0,0.4) -- (axis cs:0,5.6);
  \addplot[draw=ruleGray, line width=1.3pt, forget plot] coordinates {(33.3,1) (-20.0,1)};
  \addplot[draw=ruleGray, line width=1.3pt, forget plot] coordinates {(35.8,2) (-10.0,2)};
  \addplot[draw=ruleGray, line width=1.3pt, forget plot] coordinates {(21.0,3) (-28.2,3)};
  \addplot[draw=ruleGray, line width=1.3pt, forget plot] coordinates {(46.7,4) (10.0,4)};
  \addplot[draw=ruleGray, line width=1.3pt, forget plot] coordinates {(17.9,5) (-6.7,5)};
  \addplot[only marks, mark=o, mark size=2.6pt, thick, color=ink, fill=white,
    mark options={fill=white}]
    coordinates {(33.3,1) (35.8,2) (21.0,3) (46.7,4) (17.9,5)};
  \addlegendentry{design tasks}
  \addplot[only marks, mark=*, mark size=2.6pt, color=ours]
    coordinates {(-20.0,1) (-10.0,2) (-28.2,3) (10.0,4) (-6.7,5)};
  \addlegendentry{held-out tasks}
\end{axis}
\end{tikzpicture}
\caption{Transfer from design tasks to held-out tasks. Each lever's gain over its same-batch base in
pass rate, on the tasks it was designed on (hollow) and on held-out tasks (filled). The V1 rows
come from small sets (5 to 28 runs a side, \cref{tab:structures}); the O2 row compares the readable
tasks that split across draws (28 trials a side) with the unread 30 (two draws a side). Only the
acceptance checker stayed ahead out of sample.}
\label{fig:transfer}
\end{figure}

\begin{enumerate}[leftmargin=1.6em]
  \item \textbf{Same-model redundancy adds little information.} The misses that remain are one
    reading of an ambiguous sentence, made the same way in most draws. Readers, jurors, graders and
    checkers from the same model share that reading, so they confirm it. This is the
    self-correction limit of \citet{huang2024cannot} at the level of a team: without a new fact, a
    second copy of the model has nothing the first lacked.
  \item \textbf{Global triggers trade one failure family for another.} A check fired on every
    delivery fixes the tasks it was designed on and reworks tasks that were already right. O2 lost
    2.0 on the unread set and another global check lost 1.5, while the real-time note gained 1.0
    on the seven tasks it aimed at and lost 1.0 on the other 26.
  \item \textbf{Advice does not move judgment.} Profile text, rules and real-time notes changed
    what the model did next, not what it chose. Reading rules even shrank the work by licensing the
    model to skip what the text did not state.
  \item \textbf{Token cuts that break the cache cost more than they save.} Moving old tool results
    out at 60k tokens saved about 10k tokens a pass but raised uncached tokens by 29\% and cost by
    14\%.
  \item \textbf{Design-set wins are in-sample.} Levers were designed from readable trajectories, so
    readable screens are in-sample and overstate them; this is the adaptive overfitting that an
    unread holdout guards against \citep{dwork2015reusable}. One lever passed a task 3 of 3 times on
    the screen and 1 of 3 on the next draw.
\end{enumerate}

\para{Our own forecasts were part of the problem} We wrote a forecast before every batch. For
levers designed on readable tasks the forecasts were often wrong in sign: O2 was forecast at
$+1.5$ on the unread set and measured $-2.0$, the real-time note at $+1.0$ and measured $0.0$.
Registration did not make the forecasts better, but it made the misses visible.

\para{Implications for harness builders} Five practices follow from these results, each tied to a
measurement above.
\begin{itemize}[leftmargin=1.4em]
  \item Make the prompt prefix stable before cutting tokens, and measure cache hits per request; a
    cut that breaks the cache costs more than it saves (\cref{sec:context}).
  \item For latency, count the lead's turns, not the size of its prompts (\cref{sec:context}).
  \item Before adding agents of the same model, measure how often a task splits across draws. When
    misses are the same reading in every draw, extra copies add cost, not information
    (\cref{sec:structure,sec:judgment}).
  \item Trigger runtime help on states that passing runs never show, such as an early failure
    hand-back, rather than on every delivery (\cref{sec:structure}).
  \item Keep a held-out set unread, run every arm beside its base in the same hours, and state the
    resolution of the instrument before reading a result (\cref{sec:method}).
\end{itemize}
